\begin{lemma}[Filtered extension by zero for equivariant supports]
\label{lemma-equivariant-support-filtered-extension-by-zero}
\begin{reference}
\cite[Chapter V, \S 5.7, formulae (5.7.5) and (5.7.6),
p. 220]{Tohoku}
\end{reference}
\begin{history}
The cited passage assumes that every member of the support family on the
orbit space has a neighbourhood in that family and then asserts the derived
filtered-colimit formula (5.7.6). The neighbourhood hypothesis proves the
degree-zero formula (5.7.5), but by itself it does not prove (5.7.6).
Indeed, the filtered-colimit theorem for right derived functors produces
the right derived functors of the composites
$\mathcal A\mapsto\Gamma_X^G(\mathcal A_U)$. Identifying these with
$H^n(X;G,\mathcal A_U)$ requires the extension-by-zero acyclicity stated
explicitly below. Exactness of extension by zero alone does not supply that
acyclicity. Moreover, the ordinary supported-cohomology comparison invoked
later in the cited passage follows from Corollary 1 to Proposition 3.10.1
under the additional hypothesis that the induced family on $X$ is
paracompactifying. We keep these two requirements separate and make no
claim that either follows from the neighbourhood hypothesis.
\end{history}
Let $G$ act on a topological space $X$, let $Y=X/G$, and let
$q:X\to Y$ be the orbit map. Let $\Psi$ be a family of supports on $Y$,
put
$$
\Phi=q^\dagger\Psi,
$$
as in Lemma
\ref{lemma-orbit-space-equivariant-support-spectral-sequence}, and assume
that every $C\in\Psi$ has a closed neighbourhood belonging to $\Psi$; in
other words, there is a $B\in\Psi$ such that
$$
C\subset\mathop{\rm Int}_Y(B).
$$
Let $\mathcal U_\Psi$ be the set of open subsets $U\subset Y$ such that
$\overline U\in\Psi$, ordered by inclusion. For $U\in\mathcal U_\Psi$ put
$$
X_U=q^{-1}(U),\qquad j_U:X_U\longrightarrow X,
\qquad
E_U(\mathcal A)=\mathcal A_U=(j_U)_!(\mathcal A|_{X_U}).
$$
Then:
\begin{enumerate}
\item The set $\mathcal U_\Psi$ is directed. Each $E_U$ is an exact
endofunctor of $\textit{Mod}_G(\underline{\mathbf Z}_X)$, and for
$U\subset V$ there is a functorial monomorphism
$E_U\to E_V$.
\item For every abelian $G$-sheaf $\mathcal A$ on $X$, there is a
functorial equality
$$
\Gamma^G_\Phi(X,\mathcal A)
=
\colim_{U\in\mathcal U_\Psi}\Gamma_X^G(\mathcal A_U).
$$
\item Assume the following extension-by-zero acyclicity condition:
for every $U\in\mathcal U_\Psi$ and every injective abelian $G$-sheaf
$\mathcal I$ on $X$, the sheaf $\mathcal I_U$ is right acyclic for
$\Gamma_X^G$, or equivalently
\begin{equation}
\label{equation-equivariant-support-extension-zero-acyclic}
H^n(X;G,\mathcal I_U)=0
\quad\text{for all }n>0.
\end{equation}
Then the canonical maps
$$
\colim_{U\in\mathcal U_\Psi}H^n(X;G,\mathcal A_U)
\longrightarrow
H^n_\Phi(X;G,\mathcal A)
$$
are isomorphisms for every $n\in\mathbf Z$. In nonnegative degrees these
isomorphisms form an isomorphism of cohomological $\delta$-functors.
\item If, in addition, $\Phi$ is paracompactifying, then the canonical
maps of $G$-modules
$$
\colim_{U\in\mathcal U_\Psi}H^q(X,\mathcal A_U)
\longrightarrow
H^q_\Phi(X,\mathcal A)
$$
are isomorphisms for every $q\in\mathbf Z$, functorially in $\mathcal A$
and compatibly with connecting maps.
\end{enumerate}
\end{lemma}

\begin{proof}
The empty set belongs to $\mathcal U_\Psi$. If
$U,V\in\mathcal U_\Psi$, then
$$
\overline{U\cup V}=\overline U\cup\overline V\in\Psi,
$$
so $U\cup V\in\mathcal U_\Psi$. Thus this set is directed.
The open subset $X_U$ is $G$-invariant, and hence extension by zero gives
$\mathcal A_U$ its natural $G$-sheaf structure. Restriction and extension
by zero are exact; for the latter assertion see Lemma
\ref{lemma-extension-by-zero-locally-closed}. If $U\subset V$, the
canonical map $\mathcal A_U\to\mathcal A_V$ is a monomorphism, as is seen
on stalks. This proves (1).

The counit maps $\mathcal A_U\to\mathcal A$ give an injective map
$$
\colim_{U\in\mathcal U_\Psi}\Gamma(X,\mathcal A_U)
\longrightarrow
\Gamma(X,\mathcal A),
$$
because all transition maps are injective. A section in the source which
comes from $U$ has support $S\subset X_U$. Consequently
$$
\overline{q(S)}\subset\overline U\in\Psi,
$$
so $S\in\Phi$.

Conversely, let $s\in\Gamma_\Phi(X,\mathcal A)$, put
$S=\text{Supp}_X(s)$, and set $C=\overline{q(S)}$. Thus $C\in\Psi$.
Choose a closed neighbourhood $B\in\Psi$ of $C$ and put
$U=\mathop{\rm Int}_Y(B)$. Then
$$
C\subset U,
\qquad
\overline U\subset B,
$$
and hence $U\in\mathcal U_\Psi$ and $S\subset X_U$. The section $s$
therefore comes from a unique section of $\mathcal A_U$. We have proved a
functorial equality of $G$-modules
\begin{equation}
\label{equation-equivariant-support-sections-colimit}
\Gamma_\Phi(X,\mathcal A)
=
\colim_{U\in\mathcal U_\Psi}\Gamma(X,\mathcal A_U).
\end{equation}
The system on the right has injective transition maps. It may therefore be
viewed as a directed union of $G$-submodules. Taking invariants commutes
with such a union for an arbitrary group $G$: an element of the union is
fixed if and only if it is already fixed in any member containing it.
Taking invariants in
(\ref{equation-equivariant-support-sections-colimit}) proves (2).

For $U\in\mathcal U_\Psi$ consider the additive functor
$$
F_U:
\textit{Mod}_G(\underline{\mathbf Z}_X)\longrightarrow\textit{Ab},
\qquad
F_U(\mathcal A)=\Gamma_X^G(E_U(\mathcal A)).
$$
These functors form a directed system, and part (2) says that their
pointwise colimit is $\Gamma_\Phi^G(X,-)$. Let
$\mathcal A\to\mathcal I^\bullet$ be a nonnegative injective resolution in
$\textit{Mod}_G(\underline{\mathbf Z}_X)$. Such resolutions exist by
Injectives, Lemma
\ref{injectives-lemma-equivariant-modules-enough-injectives}.
Exactness of $E_U$ gives a resolution
$$
\mathcal A_U\longrightarrow E_U(\mathcal I^\bullet).
$$
Under hypothesis
(\ref{equation-equivariant-support-extension-zero-acyclic}) its terms are
right acyclic for $\Gamma_X^G$.
Leray's acyclicity lemma, Derived Categories, Lemma
\ref{derived-lemma-leray-acyclicity}, therefore gives
\begin{equation}
\label{equation-equivariant-support-derived-extension-zero}
R^nF_U(\mathcal A)
=
H^n(X;G,\mathcal A_U).
\end{equation}
The source category has enough injectives and abelian groups satisfy AB5.
Derived Categories, Lemma
\ref{derived-lemma-filtered-colimit-right-derived-functors}, applied to the
system $(F_U)$, now gives
$$
\begin{aligned}
\colim_{U\in\mathcal U_\Psi}H^n(X;G,\mathcal A_U)
&=\colim_{U\in\mathcal U_\Psi}R^nF_U(\mathcal A)\\
&=R^n\Gamma_\Phi^G(X,\mathcal A)\\
&=H^n_\Phi(X;G,\mathcal A).
\end{aligned}
$$
The same lemma supplies functoriality and compatibility with connecting
maps. This proves (3).

Assume finally that $\Phi$ is paracompactifying. Let $\mathcal V_\Phi$ be
the directed set of open subsets $V\subset X$ such that
$\overline V\in\Phi$. Lemma
\ref{lemma-family-supports-colimit-extension-by-zero} gives
\begin{equation}
\label{equation-equivariant-support-ordinary-all-opens-colimit}
\colim_{V\in\mathcal V_\Phi}H^q(X,\mathcal A_V)
=H^q_\Phi(X,\mathcal A).
\end{equation}
The opens $X_U$, for $U\in\mathcal U_\Psi$, form a cofinal subsystem of
$\mathcal V_\Phi$. First,
$$
q\left(\overline{q^{-1}(U)}\right)\subset\overline U,
$$
so $\overline{q^{-1}(U)}\in\Phi$. Conversely, let
$V\in\mathcal V_\Phi$ and put
$$
C=\overline{q(\overline V)}\in\Psi.
$$
Choose a closed neighbourhood $B\in\Psi$ of $C$ and put
$U=\mathop{\rm Int}_Y(B)$. Then $U\in\mathcal U_\Psi$ and
$$
V\subset q^{-1}(U).
$$
Thus cofinality changes the left hand side of
(\ref{equation-equivariant-support-ordinary-all-opens-colimit}) into the
colimit asserted in part (4).
Every $X_U$ is $G$-invariant, all transition maps are equivariant, and the
isomorphism
(\ref{equation-equivariant-support-ordinary-all-opens-colimit}) is natural
for the action isomorphisms defining the $G$-sheaf $\mathcal A$. Hence the
cofinality isomorphism is an isomorphism of $G$-modules. This proves (4).
\end{proof}


\begin{lemma}[Filtered second spectral sequence with supports]
\label{lemma-equivariant-support-filtered-second-spectral-sequence}
\begin{reference}
\cite[Chapter V, \S 5.7, formulae (5.7.7) and (5.7.7 bis),
p. 220]{Tohoku}
\end{reference}
\begin{history}
The three conditions listed after formula (5.7.7 bis) in the cited source
are three alternative sufficient conditions for commuting group cohomology
with the indicated filtered colimit; they are not simultaneous hypotheses.
They concern the identification of the $E_2$-page only and do not imply the
extension-by-zero acyclicity
(\ref{equation-equivariant-support-extension-zero-acyclic}) of Lemma
\ref{lemma-equivariant-support-filtered-extension-by-zero}. Accordingly,
the latter acyclicity remains an explicit hypothesis here.
\end{history}
In the situation of Lemma
\ref{lemma-equivariant-support-filtered-extension-by-zero}, assume its
extension-by-zero acyclicity condition
(\ref{equation-equivariant-support-extension-zero-acyclic}). For every
abelian $G$-sheaf $\mathcal A$ on $X$, there is an additive first-quadrant
spectral sequence,
functorial in $\mathcal A$ from the $E_2$-page onward, with
\begin{equation}
\label{equation-equivariant-support-filtered-second-spectral-sequence}
{}^{\mathrm {II}}E_2^{p,q}(\mathcal A)
=
\colim_{U\in\mathcal U_\Psi}
H^p\bigl(G,H^q(X,\mathcal A_U)\bigr)
\quad\Longrightarrow\quad
H^{p+q}_\Phi(X;G,\mathcal A).
\end{equation}
It induces a finite filtration on every abutment group.

Assume in addition that $\Phi$ is paracompactifying. Then the $E_2$-page
in
(\ref{equation-equivariant-support-filtered-second-spectral-sequence}) has
the usual form
\begin{equation}
\label{equation-equivariant-support-usual-second-page}
{}^{\mathrm {II}}E_2^{p,q}(\mathcal A)
=
H^p\bigl(G,H^q_\Phi(X,\mathcal A)\bigr)
\end{equation}
if any one of the following three conditions holds:
\begin{enumerate}
\item $G$ is finite;
\item for every $q\geq 0$ and $U\in\mathcal U_\Psi$, the action of $G$
on $H^q(X,\mathcal A_U)$ is trivial, and every group
$H_r(G,\mathbf Z)$ is finitely generated;
\item for every $q\geq 0$, there is a cofinal subsystem of
$\mathcal U_\Psi$ on which all transition maps
$$
H^q(X,\mathcal A_U)\longrightarrow H^q(X,\mathcal A_V)
$$
are isomorphisms.
\end{enumerate}
\end{lemma}

\begin{proof}
For each $U\in\mathcal U_\Psi$, apply the second spectral sequence of
Lemma \ref{lemma-orbit-space-equivariant-spectral-sequences} to the
abelian $G$-sheaf $\mathcal A_U$. This gives
\begin{equation}
\label{equation-equivariant-support-local-second-spectral-sequence}
{}^{\mathrm {II}}E_2^{p,q}(\mathcal A,U)
=H^p\bigl(G,H^q(X,\mathcal A_U)\bigr)
\quad\Longrightarrow\quad
H^{p+q}(X;G,\mathcal A_U).
\end{equation}
For $U\subset V$, the equivariant monomorphism
$\mathcal A_U\to\mathcal A_V$ and functoriality of that lemma give a
morphism from the spectral sequence for $U$ to the one for $V$. Hence the
spectral sequences
(\ref{equation-equivariant-support-local-second-spectral-sequence}) form a
directed system from their $E_2$-pages onward.

For $r\geq 2$, take the filtered colimit of the $r$th pages and their
differentials. Filtered colimits of abelian groups are exact by Algebra,
Lemma \ref{algebra-lemma-directed-colimit-exact}; consequently they commute
with kernels, images, and the cohomology of complexes. Therefore
$$
H\left(\colim_U{}^{\mathrm {II}}E_r(\mathcal A,U),
\colim_U d_r\right)
=
\colim_U{}^{\mathrm {II}}E_{r+1}(\mathcal A,U).
$$
These colimit pages and differentials thus form a spectral sequence, and
its $E_2$-page is the one displayed in
(\ref{equation-equivariant-support-filtered-second-spectral-sequence}).

We spell out convergence. In a first-quadrant cohomological spectral
sequence, the incoming differential at $(p,q)$ is zero for $r>p$, and the
outgoing differential is zero for $r>q+1$. Hence the terms at $(p,q)$
stabilize, uniformly in $U$, once
$$
r>\max\{p,q+1\}.
$$
It follows that the stable term of the colimit spectral sequence is the
colimit of the stable terms. For a fixed total degree $n$, let
$F_U^pH^n(X;G,\mathcal A_U)$ be the finite filtration supplied by
(\ref{equation-equivariant-support-local-second-spectral-sequence}),
and set
$$
F^p\left(\colim_UH^n(X;G,\mathcal A_U)\right)
=
\colim_UF_U^pH^n(X;G,\mathcal A_U).
$$
Exactness of filtered colimits identifies the successive quotients with
the corresponding colimits of the successive quotients, hence with
${}^{\mathrm {II}}E_\infty^{p,n-p}(\mathcal A)$. The filtration is finite,
since $0\leq p\leq n$. Part (3) of Lemma
\ref{lemma-equivariant-support-filtered-extension-by-zero} identifies its
underlying abutment with $H^n_\Phi(X;G,\mathcal A)$. This proves convergence
and the filtration assertion. Additivity and functoriality are inherited
from
(\ref{equation-equivariant-support-local-second-spectral-sequence}) and
from filtered colimits.

It remains to prove the three interchange criteria. Put
$$
M_U^q=H^q(X,\mathcal A_U).
$$
For every $p,q$ there is a canonical comparison map
\begin{equation}
\label{equation-equivariant-support-group-colimit-comparison}
\colim_UH^p(G,M_U^q)
\longrightarrow
H^p\left(G,\colim_UM_U^q\right).
\end{equation}

If $G$ is finite, its bar resolution has in every degree a finitely
generated free $\mathbf Z[G]$-module. Applying
$\Hom_{\mathbf Z[G]}(-,M)$ gives a cochain complex whose terms commute with
filtered colimits. Exactness of filtered colimits then shows that its
cohomology commutes with filtered colimits. Thus
(\ref{equation-equivariant-support-group-colimit-comparison}) is an
isomorphism in case (1).

Assume (2). For a trivial $G$-module $M$, the universal coefficient
sequence gives, naturally in $M$,
\begin{equation}
\label{equation-equivariant-support-group-universal-coefficient}
0\longrightarrow
\Ext^1_{\mathbf Z}(H_{p-1}(G,\mathbf Z),M)
\longrightarrow H^p(G,M)
\longrightarrow
\Hom_{\mathbf Z}(H_p(G,\mathbf Z),M)
\longrightarrow0,
\end{equation}
where the first term is zero for $p=0$. A finitely generated abelian group
is finitely presented, and both $\Hom_{\mathbf Z}(N,-)$ and
$\Ext^1_{\mathbf Z}(N,-)$ commute with filtered colimits when $N$ is
finitely generated. Taking filtered colimits of
(\ref{equation-equivariant-support-group-universal-coefficient}), and using
their exactness, proves that
(\ref{equation-equivariant-support-group-colimit-comparison}) is an
isomorphism in case (2).

Assume (3). On the stated cofinal subsystem every map $M_U^q\to M_V^q$
is an isomorphism. The map from any one of its terms to its colimit is
therefore an isomorphism, and the same is true after applying
$H^p(G,-)$. Thus
(\ref{equation-equivariant-support-group-colimit-comparison}) is an
isomorphism in case (3), without any condition on $G$.

Under any one of the three conditions,
(\ref{equation-equivariant-support-group-colimit-comparison}) identifies
the $E_2$-page in
(\ref{equation-equivariant-support-filtered-second-spectral-sequence}) with
$$
H^p\left(G,\colim_UM_U^q\right).
$$
Part (4) of Lemma
\ref{lemma-equivariant-support-filtered-extension-by-zero} identifies the
inner colimit, as a $G$-module, with $H^q_\Phi(X,\mathcal A)$. This gives
(\ref{equation-equivariant-support-usual-second-page}) and completes the
proof.
\end{proof}
